\section{PARC lifecycle calculus}
\label{sec:parc-lifecycle-calculus}

PARC is a lifecycle rule for finite scientific candidate queues under
one-sided partial verification.  The object is not only a selected set.  The
object is a release card: a versioned statement of the candidate universe,
reference label version, one-sided support source, score, block construction,
release/refusal decision and recertification status.

\paragraph{Setup.}
Let \(P\) be a frozen finite candidate universe and let \(Y_p\in\{0,1\}\)
denote validity under a fixed reference version.  One-sided support is a
binary observation \(A_p\) satisfying
\[
  A_p=1 \Rightarrow Y_p=1 .
\]
For a calibration block \(b\), PARC removes only verified positives from the
calibration null superset and keeps all unverified candidates in the null
superset.  Candidate scores and block definitions are frozen before release.

\begin{theorem}[Least-favourable null-superset dominance]
If one-sided reliability holds, every false calibration candidate remains in
the calibration null superset.  Therefore the maximum score over the
calibration null superset is least-favourable for the false-candidate score in
the corresponding exchangeable test block.  The resulting block rank
\(p\)-value is super-uniform for false candidates, and any monotone
calibrating transform that yields an e-value remains valid for false
candidates.
\end{theorem}

\begin{proof}
One-sided reliability removes only candidates known to be valid.  Hence no
false candidate is removed from the calibration null superset.  The block
maximum over this superset is at least as large as the maximum over the
unobserved false-only subset.  Under the stated block-exchangeability
condition, ranking a false test score against this superset maximum is
conservative.  The e-value statement follows from the standard
super-uniform-to-e-value transform used by PARC.
\end{proof}

\begin{proposition}[Refusal lower bound]
For a requested budget \(K\), risk target \(\alpha\), and compatible release
set \(R\), PARC self-consistency requires
\[
  E_p \ge \frac{K}{\alpha |R|}
  \quad \text{for all } p\in R .
\]
If no compatible non-empty set satisfies this inequality, PARC returns
certified refusal.  This refusal is a statement about evidence insufficiency
for the frozen release card, not a claim that no valid candidates exist.
\end{proposition}

\begin{proposition}[Active audit gain]
Consider a calibration candidate \(q\) that is currently a high-scoring member
of the calibration null superset.  If auditing \(q\) returns one-sided support
\(A_q=1\), then \(q\) is removed from the null superset.  This can weakly reduce
the relevant calibration block maximum and can increase the e-values of
follow-up candidates whose scores are compared against that block.  Targeted
audit policies can therefore turn certified refusal into certified release
without treating unverified candidates as negatives.
\end{proposition}

\begin{proposition}[Versioned certificate accounting]
Let \(Y_p^{(0)}\) and \(Y_p^{(1)}\) denote validity under reference versions
\(t_0\) and \(t_1\).  For any release set \(R\),
\[
 \mathrm{FTR}_{1}(R)
 =
 \mathrm{FTR}_{0}(R)
 + \frac{|\{p\in R:Y_p^{(0)}=1,Y_p^{(1)}=0\}|}{|R|\vee 1}
 - \frac{|\{p\in R:Y_p^{(0)}=0,Y_p^{(1)}=1\}|}{|R|\vee 1}.
\]
Thus inherited current-reference burden is the old-reference burden plus
stable-to-unstable drift minus unstable-to-stable correction.
\end{proposition}

\begin{theorem}[Versioned recertification]
Fix a new reference version \(t_1\).  If candidate universe, score, blocks,
compatibility and release grid are frozen, and if \(t_1\)-version one-sided
support satisfies \(A_p^{(1)}=1\Rightarrow Y_p^{(1)}=1\), then rerunning PARC
at \(t_1\) restores an expected \(t_1\)-relative false-release guarantee
whenever the returned set is non-empty and self-consistent.  If no such set is
found, PARC returns recertified refusal.
\end{theorem}

\paragraph{Lifecycle interpretation.}
The same calculus supports several release-card states: certified release,
certified refusal, active-audit requirement, expiry after reference update,
recertified release, recertified refusal and risk-triage requirement.  The
states are mutually scoped to the recorded reference version and evidence
sources.  They should not be read as prospective materials discovery, DFT
validation or absolute physical truth.
